\usepackage[T1]{fontenc}
\usepackage[utf8]{inputenc}
\usepackage[spanish,es-nodecimaldot,es-noquoting,es-lcroman]{babel}
\spanishplainpercent
\usepackage{lmodern}
\usepackage[a4paper,margin=21mm,headheight=15pt]{geometry}
\usepackage{amsmath,amssymb,bm}
\usepackage{booktabs,array}
\usepackage{tikz}
\usetikzlibrary{arrows.meta,positioning,calc}
\usepackage[most]{tcolorbox}
\usepackage{microtype}
\usepackage{fancyhdr}
\usepackage{hyperref}
\hypersetup{colorlinks=true,linkcolor=blue!45!black,urlcolor=blue!45!black}
\setlength{\parindent}{0pt}
\setlength{\parskip}{5pt}
\setlength{\emergencystretch}{2em}
\pagestyle{fancy}
\fancyhf{}
\fancyhead[L]{\small BEM del cilindro: paso \paso}
\fancyhead[R]{\small Guía matemática independiente}
\fancyfoot[C]{\thepage}
\newcommand{\vct}[1]{\bm{#1}}
\newcommand{\dd}{\,\mathrm{d}}
\newcommand{\ii}{\mathrm{i}}
\newcommand{\norm}[1]{\left\lVert#1\right\rVert}
\newcommand{\VP}{\operatorname{VP}}
\newcommand{\tr}{\operatorname{tr}}
\newtcolorbox{idea}[1]{breakable,colback=blue!3,colframe=blue!50!black,
  title={#1},fonttitle=\bfseries,boxrule=0.7pt}
\newtcolorbox{alerta}[1]{breakable,colback=orange!4,colframe=orange!65!black,
  title={#1},fonttitle=\bfseries,boxrule=0.7pt}
\newtcolorbox{resultado}[1]{breakable,colback=teal!3,colframe=teal!55!black,
  title={#1},fonttitle=\bfseries,boxrule=0.7pt}
\newcommand{\inicio}[2]{%
 {\small\color{blue!55!black}\bfseries SERIE: LOS SIETE PASOS MATEMÁTICOS DEL BEM}\par
 {\LARGE\bfseries Paso \paso. #1\par}\vspace{3mm}
 \begin{idea}{¿Qué busca este paso?}#2\end{idea}}
\newcommand{\mapa}[3]{%
 \begin{center}
 \begin{tikzpicture}[node distance=4mm,
   caja/.style={draw=blue!50!black,fill=blue!3,rounded corners=2pt,
     text width=13cm,align=center,inner sep=6pt},
   >={Stealth[length=2mm]}]
 \node[caja] (a){#1};
 \node[caja,below=of a] (b){#2};
 \node[caja,below=of b] (c){#3};
 \draw[->,thick,blue!60!black] (a)--(b);
 \draw[->,thick,blue!60!black] (b)--(c);
 \end{tikzpicture}
 \end{center}}
\newcommand{\base}{%
 \smallskip{\footnotesize\color{black!65}
 Modelo de referencia: cilindro sólido, material homogéneo e isótropo,
 deformación plana, pequeñas deformaciones, sin fuerzas de volumen;
 amplitudes complejas con \(e^{+\ii\omega t}\), \(\omega>0\).}\par}
\newcommand{\fuentes}{%
 \par\medskip\begingroup\footnotesize
 \textbf{Referencia y continuación.}
 Este documento desarrolla una etapa del esquema
 \texttt{esquema\_ecuaciones\_bem.tex}. Las ecuaciones y los ejemplos se
 contrastaron con \texttt{bem\_cilindro.py}, ambos en la carpeta superior.
 La guía general y sus deducciones ampliadas se encuentran en
 \texttt{../guia\_bem/guia\_bem.pdf}. No se modifica la implementación.
 \par Para recompilar este documento, mantener \texttt{estilo.tex} en
 la misma carpeta y ejecutar \texttt{pdflatex} sobre su archivo
 \texttt{.tex}, desde esa carpeta.\par\endgroup}
